% Statement and cited-input packet. Target proof is separately embargoed.
% Frozen II source SHA256: 74497bfe327fae25c9cdcb17334b3bf42a9793a9df7c0a4404d97568f2f760f0
% Target owner: cdp:st:D.1.6; tag: 01G4

% TARGET HELD: cdp:st:D.1.6; no separate proof environment.
% A safe statement/argument split or provenance-owner disposition is required before derivation.


% Cited input cdp:st:D.1.13; statement or labelled formula only.
\begin{proposition}[Refuted printed surjectivity and positivity, with a second example]\hypertarget{tag-01FY}{}\label{cdp:st:D.1.13}



For Wang's threefold $Y=X_2(\lambda)$ in \ref{cdp:st:D.1.6} one has
\[
K_Y\simeq f^*\OO_{\PP^1}(-1),\qquad
R^1f_*\OO_Y\simeq\OO\oplus\OO(-1),\qquad
R^2f_*\OO_Y\simeq\OO(-1).
\]
Consequently $h^q(Y,\OO_Y)=(1,1,0,0)$, and restriction to every
smooth two-torus fibre has rank one. This threefold satisfies
$H^1(Y;\Z)=H^2(Y;\Z)=0$, $a(Y)=1$, and is a second
counterexample to CDP20 Corollary 5.2 and Proposition 5.5.
Its integral homology torsion distinguishes it from the retained
integral homology sphere.
\end{proposition}

